\section{A feasible action for each world is not yet a policy}
\label{sim:certified-observation-policies}

Return to the first coupon, when cash is 5.5 and the once-only facility is available. Two companies can produce the same coarse observation that the market rate is high. In one, a fixed coupon is 5; in the other, the floating coupon at the cap is 6. Paying is feasible for the first company, while drawing is barred by the strict trigger. The second company cannot pay from existing cash and needs the facility. If the manager sees only the coarse high signal, the two individually feasible instructions do not form a decision rule the manager can follow. The rule asks for information that the manager does not have.

This is the difference between a financially possible action and an information-feasible policy. Let $w$ denote a possible history and $O_t(w)$ the manager's observation history when the date-$t$ instruction is chosen. An admissible policy must satisfy
\begin{equation}
 O_t(w)=O_t(w')\quad\Longrightarrow\quad
 \pi_t\bigl(O_t(w)\bigr)=\pi_t\bigl(O_t(w')\bigr).
 \label{eq:certified-information-policy}
\end{equation}
Here $O_t$ includes remembered earlier observations and the manager's own earlier instructions. The condition gives one decision to each information cell. It does not require that the underlying histories become identical or that the manager eventually learns the contract type.

\subsection{Refining an observation without inventing knowledge}

The selected first coarse signal reports whether $r_1\leq0.04$. The refined signal adds whether the first coupon exceeds 5.5. In the pair above, refinement distinguishes the two required first instructions. It still does not reveal the facility right or generally identify a persistent coupon type. At the second date the manager remembers the first signal and instruction, then observes whether the second coupon crosses the trigger and whether the cash then available covers it. The receipt hypothesis is held fixed for each comparison.

A financial region need not have a single observation label. For example, the condition that drawing plus cash covers a coupon is weaker than the condition that existing cash alone covers it. A region on which a draw is financially possible can therefore straddle the manager's ordinary cash-coverage observation. The checked construction splits that region on the cash-coverage inequality and its exact complement. Empty children carry the contradiction proofs of the preceding section; nonempty children retain their symbolic financial outcomes. Whole-atom proofs likewise establish that the first low/high signal and each coupon-trigger signal are constant where they are used. The observation is thus valid throughout the cell, rather than assigned from one representative rate.

This timing is consequential. The first decision does not receive the second rate or a later cash balance. The second decision uses cash before its proposed payment. Using terminal cash would let the result of an action determine what the manager supposedly observed before choosing it. Including the hidden facility right in the signal would change the information problem, even if every resulting journal still balanced.

\subsection{Which instructions survive every indistinguishable history?}

For each possible first decision rule, retain every history consistent with that rule. A first-date refusal rejects the rule. At each reachable second-date information cell, intersect the allowed instructions across the retained histories. An empty intersection means there is no action that works in every history the manager still considers possible. Otherwise choose one instruction from each cell's intersection. Cells that cannot be reached under the first rule need no arbitrary decision in this comparison.

An independent formulation expresses the same condition as Boolean restrictions. Let $X_o$ say whether the first instruction at signal $o$ is draw, and $Y_h$ the second instruction at remembered history $h$. A certified first-date refusal under instruction $a$ forbids $X_o=a$. A certified second-date refusal under first instruction $a$ and second instruction $b$ forbids their joint selection:
\begin{equation}
 (X_o\ne a)\ \lor\ (Y_h\ne b).
\end{equation}
Enumerating the satisfying assignments and retaining only reachable second-date choices gives the complete policy set. This independently authored constraint calculation agrees with the action-intersection calculation and with the earlier finite-policy reference. Agreement concerns whole policies, not only the number found.

Across the eight supplied problems---receipts 6 or 7, coarse or refined observations, and committed or mixed facility rights---only the receipt-7, refined-observation, committed-right problem has robustly feasible policies. It has four. The first instruction is fixed by the refined trigger. After a first draw, the exhausted facility cannot be used again and the second instruction must be pay. After a first pay, two reachable second-date cells allow either pay or draw, giving $2\times2=4$ reachable policies. The model supplies no preference or objective that ranks them.

The failures have different causes. Coarse information combines incompatible first instructions. A family containing an unavailable facility includes a high-coupon world with no feasible first instruction. Receipt 6 leaves the double-high path short after repayment, despite knowing the trigger and having the original right. Removing such a world or ignoring its refused branch changes the robust question; it does not solve it.

These results concern the explicit family of possible histories and the supplied observation maps. They assign no probabilities to that family, and a feasible robust policy is not a forecast or an optimal policy. Applying the calculation to an actual manager next requires evidence about source revisions, access times and observations. That evidence must preserve what was available when a decision was made rather than replacing it with a later corrected account.\footnote{The complete observation refinements, whole-atom proofs, two policy formulations, source locks and eight complete comparisons are retained in \path{data/accounting-master/certified-observation-policy-20261026/attempt-01/}. The reference checks all 3,872 saved financial jobs, with 1,766 complementary observation refinements. These are exact calculations for the selected synthetic information structure.}
